% Глава: 10-tmux
% Статус: наполнено (первый прогон).
\chapter{Управление сессиями: tmux}
\label{ch:10-tmux}

\section{Зачем нужен tmux}

При работе по SSH (\nameref{ch:02-linux-access}) разрыв
соединения убивает все запущенные в терминале процессы.
\textbf{tmux} --- менеджер терминальных сессий: он позволяет
запустить программу в \emph{сессии}, отключиться, вернуться
позже (с другого компьютера) и продолжить работу --- процессы
останутся.

Для работы на сервере \fpath{pglab.nsu.ru} это особенно важно:
длинная компиляция или тесты не прервутся, даже если вы
случайно закрыли SSH-сессию.

Основные понятия:
\begin{center}
\begin{tabular}{lp{0.7\linewidth}}
\toprule
\textbf{Понятие} & \textbf{Что это} \\
\midrule
session & именованное пространство (несколько окон) \\
window  & «вкладка» внутри сессии \\
pane    & область внутри окна (деление экрана) \\
\bottomrule
\end{tabular}
\end{center}

\begin{figure}[htb]
  \centering
  \begin{tikzpicture}[node distance=1.2cm]
    \node[process, minimum width=6cm] (s) {session: \texttt{work}};
    \node[ext, below left=1.1cm and -0.5cm of s] (w1) {window: monitor};
    \node[ext, below right=1.1cm and -0.5cm of s] (w2) {window: editor};
    \node[data, below left=1.1cm and -0.3cm of w1] (p1) {pane: generator};
    \node[data, below right=1.1cm and -0.3cm of w1] (p2) {pane: tail -f};
    \node[data, below=1.1cm of w2] (p3) {pane: vim};
    \draw[flow] (s) -- (w1);
    \draw[flow] (s) -- (w2);
    \draw[flow] (w1) -- (p1);
    \draw[flow] (w1) -- (p2);
    \draw[flow] (w2) -- (p3);
  \end{tikzpicture}
  \caption{Иерархия tmux: сессия содержит окна, окно --- панели.}
  \label{fig:tmux}
\end{figure}

\section{Основное управление}

Большинство команд начинаются с \textbf{prefix} (по умолчанию
\cmd{Ctrl+b}), затем нажимается клавиша-команда. Ниже ---
без prefix там, где он не нужен.

\begin{minted}[bgcolor=gray!6]{bash}
tmux                      # новая сессия по умолчанию
tmux new -s work          # именованная сессия
tmux ls                   # список сессий
tmux attach -t work       # присоединиться к сессии
tmux kill-session -t work # удалить сессию
\end{minted}

Отключение от сессии (процессы продолжают работать):
нажмите \cmd{Ctrl+b}, отпустите, затем \cmd{d} (detach).

\section{Окна и панели (после prefix \cmd{Ctrl+b})}

\begin{center}
\begin{tabular}{lp{0.7\linewidth}}
\toprule
\textbf{Клавиша} & \textbf{Действие} \\
\midrule
\cmd{c}         & создать новое окно \\
\cmd{n} / \cmd{p} & следующее / предыдущее окно \\
\cmd{0}--\cmd{9} & перейти к окну по номеру \\
\cmd{\%}        & разделить окно вертикально \\
\cmd{\textbackslash} & разделить окно горизонтально \\
\cmd{h/j/k/l}  & переключение между панелями \\
\cmd{z}         & zoom: на весь экран / обратно \\
\cmd{x}         & закрыть текущую панель \\
\cmd{d}         & отключиться от сессии (detach) \\
\bottomrule
\end{tabular}
\end{center}

Типовая рабочая раскладка: в одной панели запущена долгая
программа, в другой --- наблюдение за логом (\cmd{tail -f}),
в третьем окне --- редактор.

\section{Настройка}

tmux читает \fpath{~/.tmux.conf}. Минимальный
полезный конфиг:
\begin{minted}[bgcolor=gray!6]{bash}
# использовать C-a как префикс вместо C-b
set -g prefix C-a
unbind C-b
bind C-a send-prefix
# нумеровать окна и панели с 1
set -g base-index 1
set -g pane-base-index 1
\end{minted}

\section{Типичные ошибки}
\begin{itemize}
  \item Нажали \cmd{Ctrl+b}, отпустили и ничего не произошло ---
    нужно нажать \emph{вторую} клавишу команды.
  \item Отключились и забыли имя сессии --- \cmd{tmux ls}.
  \item Сессия «умерла» --- её закрыли (kill) или сервер tmux
    завершился; detach от этого отличается.
  \item Не можете найти панель --- \cmd{Ctrl+b z} переключает
    zoom, \cmd{Ctrl+b q} показывает номера панелей.
\end{itemize}

\section{Связь с остальным курсом}

tmux --- стандарт для работы на сервере
(\nameref{ch:02-linux-access}) и в Docker
(\nameref{ch:19-docker-basics}): внутри контейнера можно
поднять tmux и иметь несколько панелей для наблюдения
за сервисами. Наблюдение за процессами в панелях ---
переход к \nameref{ch:12-processes}.

\section{Задачи для самопроверки}
\begin{enumerate}
  \item Чем detach от сессии отличается от kill сессии?
  \item Запустите сессию, создайте два окна и разделите одно
    окно на две панели. Как между ними переключаться?
  \item Как отключиться от сессии, не убивая процессы?
  \item Как переименовать или перейти к окну по номеру?
  \item Как вернуть панель из zoom?
\end{enumerate}
